\section{Introduzione}

Bluesky e' un social network costruito su ATProto, un protocollo che separa funzioni di identita', repository personali, distribuzione dei dati e servizi applicativi. Questa architettura rende osservabile una parte rilevante delle interazioni pubbliche e dei metadati di moderazione: post, follows, blocchi individuali, liste, labeler e labels possono essere raccolti attraverso stream o API pubbliche, entro i limiti tecnici e normativi del protocollo \parencite{blueskyAtprotoDocs, atprotoOverview, balduf2024looking}. Per la ricerca empirica sulla moderazione, questa apertura e' importante perche' consente di osservare non solo decisioni ufficiali della piattaforma, ma anche segnali prodotti dagli utenti e da servizi indipendenti.

Il progetto studia il rapporto tra moderazione comunitaria e moderazione ufficiale. La moderazione comunitaria comprende azioni individuali come i blocchi, strumenti aggregati come modlist o blocklist, e labels prodotte da labeler di terze parti. La moderazione ufficiale e' osservata tramite eventi account-level \code{!takedown} emessi dal servizio di moderazione ufficiale. Nel report, il takedown e' trattato come proxy operativo dell'enforcement ufficiale a livello account, non come verita' assoluta sullo stato o sulla causa della sospensione.

La domanda scientifica e' circoscritta: i segnali comunitari pubblicamente osservabili prima dell'evento contengono informazione sul successivo takedown account-level? La formulazione e' deliberatamente prudente. Lo studio e' osservazionale, temporalmente ancorato all'evento, correlazionale e predittivo solo in senso statistico. Non identifica effetti causali: i blocchi possono precedere temporalmente il takedown ed essere associati ad esso, ma questo non significa che causino la decisione ufficiale.

La letteratura su Bluesky ha gia' documentato la struttura modulare della piattaforma, la scala dei dati osservabili e il ruolo del blocking nella self-moderation \parencite{balduf2024looking, bono2026selfmoderation, sokoto2026open}. Tuttavia, il legame tra segnali comunitari pre-evento e successivi takedown ufficiali account-level resta meno esplorato. Il contributo di questo lavoro e' quindi empirico e metodologico: costruisce una finestra temporale pre-evento; distingue account positivi e controlli negativi con pseudo-event time; controlla per esposizione tramite bucket di post e follow ricevuti; separa la popolazione 00 dagli account esposti; confronta blocchi, labels e modlist; valuta una Random Forest basata sui soli segnali di blocco ricevuti.

I risultati anticipano una distinzione centrale. Per account con una certa esposizione pubblica, i blocchi ricevuti costituiscono un segnale informativo, concentrato soprattutto nel giorno immediatamente precedente al takedown. Per gli account nel corner 00, cioe' con zero post e zero follow ricevuti nella finestra osservata, il segnale comunitario e' invece molto debole in valore assoluto. Questa popolazione rappresenta il 60,23\% dei positivi e presenta spesso una vita molto breve, dinamica coerente con account usa-e-getta, spam, ban evasion, impersonation o altri pattern anti-abuso account-level. Il punto non e' che tutti questi account siano classificabili come bot, ma che molti takedown sembrano avvenire prima che i segnali comunitari pubblici possano formarsi.

\Cref{fig:architecture} sintetizza il quadro concettuale: ATProto rende osservabili diversi segnali pubblici; il progetto li ancora temporalmente a un evento ufficiale; il confronto empirico valuta presenza, intensita', temporalita' e capacita' predittiva dei segnali.

\begin{figure}[htbp]
  \centering
  \includegraphics[width=0.95\linewidth]{fig01_architecture.pdf}
  \caption{Architettura concettuale dello studio. I segnali comunitari osservabili vengono misurati nella finestra pre-evento e confrontati con il successivo takedown account-level ufficiale.}
  \label{fig:architecture}
\end{figure}
